% Created 2026-03-09 Mon 14:48
% Intended LaTeX compiler: pdflatex
\documentclass[11pt, letterpaper]{article}
\usepackage[utf8]{inputenc}
\usepackage[T1]{fontenc}
\usepackage{graphicx}
\usepackage{longtable}
\usepackage{wrapfig}
\usepackage{rotating}
\usepackage[normalem]{ulem}
\usepackage{amsmath}
\usepackage{amssymb}
\usepackage{capt-of}
\usepackage{hyperref}
\usepackage[margin=2.5cm]{geometry}
\usepackage[utf8]{inputenc}
\usepackage{palatino}
\usepackage{xcolor}
\author{Equipo Alpine White}
\date{\textit{{[}2026-02-16 Mon]}}
\title{Bitácora Semanal 5: Administración de Sistemas Unix}
\hypersetup{
 pdfauthor={Equipo Alpine White},
 pdftitle={Bitácora Semanal 5: Administración de Sistemas Unix},
 pdfkeywords={},
 pdfsubject={},
 pdfcreator={Emacs 30.2 (Org mode 9.7.39)}, 
 pdflang={English}}
\begin{document}

\maketitle
\setcounter{tocdepth}{2}
\tableofcontents

\section{Información General}
\label{sec:org86b3dca}

\begin{itemize}
\item \textbf{Semana:} Del 02 de Marzo 2026 al 06 de Marzo 2026
\item \textbf{Equipo:}
\begin{itemize}
\item Arreguín Salgado Gael Emiliano
\item Gramer Muñoz Omar Fernando
\item López Pérez Mariana
\item Nieto Gallegos Isaac Julián
\end{itemize}
\end{itemize}

Durante la quinta semana del curso se profundizó en diversos conceptos
relacionados con la administración de sistemas Unix, enfocándose
principalmente en redes, direcciones MAC y el uso de herramientas de
terminal para la manipulación de archivos y gestión del sistema. A lo
largo de las sesiones se reforzó la importancia del trabajo desde la
línea de comandos como medio fundamental para comprender el
funcionamiento interno de Linux, así como la interacción directa con
el hardware y los servicios del sistema operativo.
\subsection{Responsabilidades}
\label{sec:org0e03621}

Cada integrante del equipo participó activamente en el desarrollo de la
bitácora. Gael Emiliano se encargó de la recopilación de notas físicas
y del registro de comandos utilizados durante las prácticas. Omar
Fernando realizó investigaciones relacionadas con conceptos de redes y
comunicación. Mariana organizó la estructura general del documento y
su redacción principal, mientras que Isaac Julián llevó a cabo pruebas
prácticas en el sistema y documentó el comportamiento observado en los
comandos ejecutados.
\section{Direcciones MAC y redes}
\label{sec:orge8c8588}

\subsection{Conceptos nuevos}
\label{sec:orgff1fe63}

\subsubsection{Dirección MAC}
\label{sec:orgb263953}

La dirección MAC (Media Access Control) constituye un identificador
único asignado físicamente a cada interfaz de red. Este identificador
opera en la capa de enlace de datos del modelo OSI y permite distinguir
de manera inequívoca a cada dispositivo dentro de una red local. A
diferencia de la dirección IP, que puede modificarse dependiendo de la
configuración de red, la dirección MAC permanece asociada al hardware,
lo que la convierte en un elemento clave para el control y la gestión
del tráfico de red.

Su uso resulta fundamental en mecanismos como el filtrado de acceso en
routers, la identificación de dispositivos conectados y la
administración de políticas de seguridad dentro de redes locales.
\subsubsection{Formato MAC de 48 bits}
\label{sec:org1916b99}

El formato más común de dirección MAC está compuesto por seis grupos
hexadecimales separados por dos puntos. Los primeros 24 bits corresponden
al identificador del fabricante (OUI), asignado por organismos
internacionales, mientras que los bits restantes identifican de forma
única al dispositivo producido por dicho fabricante.
\subsubsection{Formato MAC de 64 bits}
\label{sec:orgb18c4f5}

En redes más modernas se emplea el formato de 64 bits, el cual amplía
el espacio de direccionamiento permitiendo soportar una mayor cantidad
de dispositivos, especialmente en entornos inalámbricos y tecnologías
recientes.
\subsubsection{ip addr}
\label{sec:orge8c074c}

El comando \texttt{ip addr} permite visualizar todas las interfaces de red del
sistema junto con información detallada como direcciones IP,
direcciones MAC y el estado operativo de cada interfaz. Esta
herramienta resulta esencial para tareas de diagnóstico, ya que
facilita verificar la conectividad y detectar configuraciones
incorrectas.

\begin{verbatim}
ip addr
\end{verbatim}
\subsubsection{SSH}
\label{sec:org33d6782}

SSH (Secure Shell) es un protocolo que permite establecer conexiones
remotas seguras entre equipos mediante cifrado. Gracias a este
mecanismo es posible administrar servidores, ejecutar comandos a
distancia y transferir información sin exponer datos sensibles a la
red.

\begin{verbatim}
ssh usuario@192.168.25.175
\end{verbatim}
\subsection{Conceptos de repaso}
\label{sec:org1f79a0a}

Se revisaron dispositivos fundamentales de red como el router, encargado
de interconectar distintas redes, y el access point, cuya función es
extender la conectividad inalámbrica. Asimismo, se analizaron los
modos de comunicación semiduplex, duplex y full duplex, diferenciados
por la forma en que los dispositivos transmiten y reciben información.
\subsection{Máximas}
\label{sec:orgd6f7e0a}

La identificación única de cada interfaz de red garantiza el correcto
funcionamiento de las comunicaciones. Asimismo, el uso de SSH representa
una evolución hacia métodos seguros de administración remota,
reemplazando protocolos inseguros utilizados anteriormente.
\subsection{Sección libre}
\label{sec:org6d093ca}

Se discutió cómo el cableado de cobre presenta pérdidas de señal
aproximadamente después de los cien metros, fenómeno causado por la
atenuación eléctrica del medio físico, lo cual influye directamente en
el diseño de infraestructuras de red.
\section{Manipulación de archivos y texto}
\label{sec:org0d94ec0}

\subsection{Conceptos nuevos}
\label{sec:org8ce17e5}

Durante esta sección se exploraron herramientas clásicas de Unix
orientadas al procesamiento eficiente de texto desde la terminal,
demostrando la filosofía del sistema basada en combinar programas
pequeños para resolver tareas complejas.
\subsubsection{tail}
\label{sec:org28e6c24}

El comando \texttt{tail} permite visualizar las últimas líneas de un archivo,
siendo especialmente útil para monitorear archivos de registro en
tiempo real sin necesidad de abrirlos completamente.

\begin{verbatim}
tail nombres
\end{verbatim}

También admite opciones para limitar el número de líneas o bytes
mostrados, facilitando la inspección rápida de información reciente.

\begin{verbatim}
tail -2 nombres
tail -c 3 nombres
\end{verbatim}
\subsubsection{head}
\label{sec:orgfe3cdcf}

A diferencia de tail, el comando \texttt{head} muestra el inicio de un archivo,
lo cual resulta práctico al analizar archivos extensos o verificar su
estructura inicial.

\begin{verbatim}
head nombres
head -3 nombres
\end{verbatim}
\subsubsection{tr}
\label{sec:orgd2840ad}

La herramienta \texttt{tr} transforma caracteres provenientes de la entrada
estándar. Permite convertir texto, eliminar caracteres o reemplazarlos
siguiendo patrones definidos, siendo ampliamente utilizada en procesos
de limpieza y normalización de datos.

\begin{verbatim}
cat nombresD | tr 'a-z' 'A-Z'
\end{verbatim}
\subsubsection{uniq}
\label{sec:orge931486}

El comando \texttt{uniq} elimina líneas repetidas consecutivas dentro de un
archivo. Generalmente se combina con \texttt{sort} para agrupar previamente
los datos y detectar duplicados de manera efectiva.

\begin{verbatim}
sort nombres | uniq
\end{verbatim}
\subsubsection{grep}
\label{sec:orgb2d7c5a}

\texttt{grep} permite buscar patrones dentro de texto mediante expresiones
regulares. Es una herramienta fundamental para filtrar información,
analizar registros del sistema y localizar datos específicos.

\begin{verbatim}
cat nombres | grep -E Er.*
\end{verbatim}
\subsubsection{sed}
\label{sec:org9660952}

El editor de flujo \texttt{sed} posibilita modificar texto automáticamente sin
abrir archivos manualmente. Puede reemplazar palabras, eliminar líneas
o transformar contenido de forma masiva, lo cual lo convierte en una
herramienta esencial para automatización.

\begin{verbatim}
sed -i 's/Erick/Ayudante/' nombres
sed '2d' nombres1
\end{verbatim}
\subsection{Conceptos de repaso}
\label{sec:org37a0312}

Se reforzó el uso de \textbf{pipes} y redirecciones, mecanismos que permiten
encadenar comandos y construir flujos de procesamiento eficientes
directamente desde la terminal.
\subsection{Máximas}
\label{sec:orgc6ac54d}

La filosofía Unix demuestra que herramientas simples, correctamente
combinadas, pueden sustituir aplicaciones complejas y aumentar la
eficiencia del trabajo.
\subsection{Sección libre}
\label{sec:orge98094c}

Se concluyó que el dominio del procesamiento textual constituye una de
las habilidades más importantes para la automatización dentro de Linux.
\section{Administración del sistema y discos}
\label{sec:orgf277edf}

\subsection{Conceptos nuevos}
\label{sec:org185cdf7}

\subsubsection{fdisk}
\label{sec:orgd95c165}

El comando \texttt{fdisk} permite administrar particiones de disco desde la
terminal. A través de esta herramienta es posible visualizar discos
disponibles, crear nuevas particiones o eliminarlas, lo cual impacta
directamente en la organización del almacenamiento.

\begin{verbatim}
fdisk -l
\end{verbatim}
\subsubsection{RAID}
\label{sec:org44a58e7}

RAID consiste en la combinación lógica de múltiples discos físicos con
el objetivo de mejorar el rendimiento o aumentar la tolerancia a fallos.
Dependiendo del nivel utilizado, puede priorizar velocidad,
redundancia o un equilibrio entre ambas características.
\subsubsection{NVMe}
\label{sec:org576c06e}

La tecnología NVMe representa una evolución en almacenamiento sólido,
ya que utiliza el bus PCIe en lugar de SATA, ofreciendo velocidades de
lectura y escritura considerablemente superiores y menor latencia.
\subsection{Conceptos de repaso}
\label{sec:orgd5b1388}

Se revisaron conceptos relacionados con discos tradicionales y la
configuración física del hardware, comprendiendo cómo estos elementos
influyen en el rendimiento general del sistema.
\subsection{Máximas}
\label{sec:org3dac536}

El rendimiento del sistema operativo se encuentra estrechamente ligado
a la velocidad y correcta configuración del almacenamiento.
\subsection{Sección libre}
\label{sec:org341db17}

Se reflexionó sobre la relación entre hardware y software, destacando
cómo decisiones físicas pueden impactar directamente la experiencia del
usuario y la eficiencia del sistema.
\section{Administración de paquetes y documentación}
\label{sec:org3627ba3}

\subsection{Conceptos nuevos}
\label{sec:org433a275}

\subsubsection{dpkg}
\label{sec:orgaa972e0}

El gestor de paquetes \texttt{dpkg} permite instalar, eliminar y consultar
software en sistemas basados en Debian. Proporciona control directo
sobre los paquetes del sistema, facilitando la administración manual
del software instalado.

\begin{verbatim}
dpkg -l
dpkg -i archivo.deb
dpkg -P paquete
\end{verbatim}
\subsubsection{chown}
\label{sec:orgd4ace4b}

El comando \texttt{chown} cambia el propietario de archivos y directorios,
permitiendo administrar permisos y controlar el acceso dentro del
sistema.

\begin{verbatim}
sudo chown -R usuario carpeta
\end{verbatim}
\subsubsection{man}
\label{sec:orgc14c289}

La herramienta \texttt{man} proporciona acceso a la documentación oficial de
cada comando instalado en el sistema, convirtiéndose en una fuente de
consulta indispensable para el aprendizaje autónomo.

\begin{verbatim}
man ls
\end{verbatim}
\subsubsection{dmidecode}
\label{sec:orgb6cb9de}

El comando \texttt{dmidecode} permite obtener información detallada del
hardware directamente desde la BIOS o UEFI, incluyendo fabricante,
modelo del equipo y características de memoria.

\begin{verbatim}
dmidecode
\end{verbatim}
\subsection{Conceptos de repaso}
\label{sec:org3ad4ae2}

Se reforzó la importancia de la documentación integrada y la correcta
gestión del software instalado.
\subsection{Máximas}
\label{sec:org254ca46}

Un administrador de sistemas eficiente utiliza primero la documentación
local antes de recurrir a fuentes externas.
\subsection{Sección libre}
\label{sec:org216efa4}

El uso constante del comando \texttt{man} fomenta la autonomía técnica y el
aprendizaje continuo dentro del entorno Linux.
\end{document}
